% ==================================================================
\section{问题重述}
% ==================================================================

\subsection{问题背景}

人表达情绪时，说了什么、用什么语气说、说话时脸上是什么表情，三者往往同时起作用。多模态情感分析就是把这三路信号放在一起判断，它是人机对话、具身智能等应用的基础环节。CMU-MOSEI 是这一方向最常用的英文视频数据集之一\cite{zadeh2018mosei}。

可是到了真实视频上，事情并不整齐。三路信号的时间尺度不同：文本以词为单位，声音以毫秒为单位分帧，画面以固定帧率抽取；片段里还常有静音、人脸出画、转写与声音对不上的情况。这带来两个直接的问题：模型在某一路局部缺失时是否还靠得住，以及模型给出判断后，能不能指出它依据的是哪句话、哪段声音、哪几帧画面。

赛题围绕这两个问题设置了三项任务：从原始视频提取特征并做时序对齐，在模态局部缺失下做鲁棒预测，以及做可解释的预测。本文依次完成这三项任务。

\subsection{数据说明与使用边界}\label{sec:data}

赛题提供四个附件，本文对它们的用法见表\ref{tab:data}。附件2的对齐版特征 aligned\_50.pkl 共4850条，训练集、验证集、测试集分别为3395、728、727条。

\begin{table}[htbp]
\centering
\caption{题目附件及本文的使用方式}
\label{tab:data}
\begin{tabularx}{\textwidth}{l>{\raggedright\arraybackslash}p{3.3cm}c>{\raggedright\arraybackslash}X}
\toprule
数据 & 规模与结构 & 标签 & 本文用途 \\
\midrule
附件1 & 100段原始视频，37个源视频 & 有 & 只用于问题一的原始特征、时间映射与异常分析，不进入问题二的训练 \\
附件2 & 训练3395、验证728、测试727；50个位置 & 有 & 训练集拟合参数，验证集选结构和检查点；测试集只在模型冻结后评估一次 \\
附件3 & 30条对齐特征 & 无 & 问题二的缺失专项推理 \\
附件4 & 20条对齐特征及原始视频 & 无 & 问题三的预测、解释与人工时间核验 \\
\bottomrule
\end{tabularx}
\end{table}

有两点需要事先说明。第一，测试集在模型冻结后我们评估过一次，已经看过结果，所以后面的改进实验一律不再使用它，也不把它当作从未接触过的盲测集。第二，问题一自己提取的声学特征是25维，附件2的声学特征是74维，两者都来自声音，却不是同一套特征；视觉虽然都是35维，含义也不一定相同。问题二、三只用附件2的现成特征，不与问题一的特征混用。

\subsection{需要解决的问题}

\textbf{问题一：}基于附件1的原始视频，建立文本、语音、视觉三模态情感特征提取与时序对齐的数学模型，说明原始样本的处理流程、各模态特征的定义与提取方法、跨模态对齐的规则与实现；生成的特征要做到样本覆盖完整、时序组织可核验、方法可复现，并用典型样本展示对齐效果。

\textbf{问题二：}针对文本、语音或视觉中部分连续时间段不可用的情形，构建鲁棒的情感预测模型，同时输出极性和强度；分析缺失模态类型、缺失位置、缺失时长对性能的影响；对附件3完成推理。

\textbf{问题三：}针对三模态完整但决策依据难以追溯的情形，构建可解释的情感预测模型，在输出极性与强度的同时，定量给出各模态的作用程度、主要参考模态，以及能对应到原始文本片段、语音时段或视觉关键帧的关键证据；对附件4完成推理。

% ==================================================================
\section{问题分析}
% ==================================================================

\subsection{问题一的分析}

问题一的难点在于让每个特征位置都能回到原始素材上。具体有三层。

第一层是时间从哪里来。文本本身没有时间戳。最省事的办法是把整段视频均匀切成若干份，按顺序分给各个词，但这样会给停顿、静音、人不在画面里的时段凭空编出“词时间”。我们更希望每个词的时间都有依据。

第二层是一词多位。BERT 按子词切分，一个英文词可能占两三个位置，这些位置应当共享同一个词的时间窗，而不是各自再切一段。

第三层是提取失败。有的片段没有语音，有的转写和声音对不上，有的全程找不到清晰的人脸。这些情况如果都用零填充，下游就分不清“这里没有表情”和“这里根本没提取到”。

据此，问题一的路线是：先用强制对齐工具 WhisperX\cite{bain2023whisperx} 在给定转写上求词级时间，再用字符重叠把子词对到原词，然后把声学帧和人脸动作帧按词时间窗池化到50个位置；凡是拿不到时间或拿不到帧的位置，都显式记录其状态，而不是悄悄补零。

\subsection{问题二的分析}

与问题一不同，问题二直接用附件2现成的50位置特征，重点从“特征怎么来”转到“缺了一段还能不能判断准”。难点有两层。

第一层是零值有三种来源：序列本身较短留下的补齐填充、原始数据里就没有观测、我们为做实验人为遮掉的部分。三者如果混在一起，模型和评估都会被误导。因此我们对每个位置分别记录这三种掩码。

第二层是评估的独立性。同一条样本会在很多个缺失档位里重复出现，同一个源视频也会切出多个片段，它们都不是彼此独立的样本。所以我们对63个缺失档先分别算指标再等权平均，比较两个模型时按源视频成组做配对重抽样。

在模型选择上，张量融合 TFN\cite{zadeh2017tfn}、低秩融合 LMF\cite{liu2018lmf} 和跨模态 Transformer MulT\cite{tsai2019mult} 代表了表达能力逐步增强的几类融合方式。但本题训练样本只有3395条，输入又已经按50个位置对齐，复杂模型不一定占优。我们因此以参数很少的池化双头模型为主线，把低秩融合、文本 BiLSTM、跨模态时序网络等作为同一评估口径下的对照，让数据来决定用哪个。

\subsection{问题三的分析}

与问题二相比，问题三的输入是完整的，难点从“预测得稳”转到“说得清”。这里要回答三个问题：模态重要性的排序针对的是哪一个输出；被选出来的词或帧，删掉以后是不是真的会改变这个输出；给出的时间位置在原始视频里是否真有其事。

这三个问题需要三种不同的检验。模态层面，我们用 Shapley 值\cite{lundberg2017shap}把完整预测与各模态子集预测之间的差分摊给三个模态，三模态只有8个子集，可以精确计算，不需要采样近似。局部证据层面，参照 ERASER\cite{deyoung2020eraser}的做法，用删除高分位置后预测的变化来检验证据是否忠实于模型。时间层面，则由人工对照视频逐词核对。前两项检验的是“模型是不是这么用的”，第三项检验的是“素材里是不是真有这个”，谁也替代不了谁。

\subsection{求解思路}

三问的整体流程如图\ref{fig:flow}所示。

\begin{figure}[htbp]
\centering
\begin{tikzpicture}[
  node distance=3mm,
  box/.style={draw=black!65, rounded corners=2pt, align=flush center, font=\small, inner sep=3pt, minimum height=7.5mm},
  wide/.style={box, text width=13.2cm, fill=black!4},
  q1/.style={box, text width=13.2cm, fill=green!6},
  q2/.style={box, text width=6.2cm, fill=blue!5},
  q3/.style={box, text width=6.2cm, fill=orange!7},
  head/.style={font=\small\bfseries, text width=6.2cm, align=flush center},
  arr/.style={-{Stealth[length=1.8mm]}, draw=black!70, semithick}]
\node[head, text width=13.2cm] (h1) {问题一：原始视频的三模态特征与时间对齐（附件1）};
\node[q1, below=2mm of h1] (p1) {WhisperX 强制对齐给定转写 $\to$ 词级时间；字符重叠 $\to$ 子词对原词};
\node[q1, below=of p1] (p2) {BERT（768维）、eGeMAPS 声学（25维）、OpenFace 动作单元（35维）按词时间窗池化到50个位置};
\node[q1, below=of p2] (p3) {可观测掩码与失败状态、人工核验、100条台账};
\node[wide, below=8mm of p3] (d) {附件2：分离补齐填充、原生不可观测、人为遮蔽三种掩码；训练集拟合、验证集选型};
\node[head, below=5mm of d.south west, anchor=north west] (h2) {问题二：局部缺失下的预测};
\node[head, below=5mm of d.south east, anchor=north east] (h3) {问题三：模态归因与局部证据};
\node[q2, below=2mm of h2] (a1) {PM-B2 池化双头模型与7个基线};
\node[q2, below=of a1] (a2) {63档局部缺失面板、视频级配对重抽样};
\node[q2, below=of a2] (a3) {缺失类型、比例、位置的影响规律};
\node[q2, below=of a3] (a4) {A–E 改进尝试与整模态压力面板};
\node[q2, below=of a4] (a5) {测试集误差分析、附件3推理};
\node[q3, below=2mm of h3] (c1) {固定目标类别，八子集精确 Shapley};
\node[q3, below=of c1] (c2) {逐位置删除：高分 / 低分 / 随机对照};
\node[q3, below=of c2] (c3) {子词位置回指原词与音视频帧};
\node[q3, below=of c3] (c4) {人工时间核验（拒绝不可信时间）};
\node[q3, below=of c4] (c5) {附件4预测与解释卡};
\draw[arr] (p1) -- (p2); \draw[arr] (p2) -- (p3);
\draw[arr, dashed] (p3) -- node[right, font=\scriptsize]{特征口径不同，只共享处理思路} (d);
\draw[arr] (d.south -| h2.north) -- (h2.north);
\draw[arr] (d.south -| h3.north) -- (h3.north);
\foreach \x/\y in {a1/a2,a2/a3,a3/a4,a4/a5,c1/c2,c2/c3,c3/c4,c4/c5} \draw[arr] (\x) -- (\y);
\draw[arr, dashed] (a1.east) -- (c1.west);
\end{tikzpicture}
\caption{全文求解思路（中间的虚线箭头表示问题三直接使用问题二冻结的正式模型）}
\label{fig:flow}
\end{figure}

% ==================================================================
\section{模型假设}
% ==================================================================

\begin{enumerate}[label=(\arabic*)]
\item 情感标签对应整个片段，片段内的50个位置共享同一个标签；位置级证据说明模型依赖哪里，不代表该词有单独的情感真值。
\item 同一源视频切出的多个片段之间可能相关，比较模型差异时按源视频成组抽样。
\item 人为遮蔽只改变观察到的特征，不改变标签；全部模态都不可用时，模型返回训练集的类别先验和强度中位数。
\item 转写与声音不符、静音或无人脸时，保留样本及其可用模态，但不为它编造词时间或视频帧。
\item 附件3、4的标签未知，其结果只能检查格式、内部一致性和人工可核对性，不能计算监督意义上的准确率。
\end{enumerate}

\textbf{说明。}假设(2)直接影响区间估计：附件2验证集的728条来自239个源视频，同一视频的片段在说话人、场景上高度相似。如果按条独立重抽样，会把这种相关性当成独立信息，得到过窄的置信区间，所以我们按视频成组抽样。

假设(3)界定了缺失面板回答的是什么问题：在预先设定的遮蔽方式下模型还剩多少能力。真实世界的缺失未必是这样发生的，因此63档结果代表的是这类可控干预下的鲁棒性，不能直接推广到任意缺失机制。

假设(4)在问题一和问题三中都用到了。宁可让一个位置“没有时间”，也不给它一个看起来合理却无法核对的时间，因为后者一旦被当作证据引用，比缺失更难发现。

% ==================================================================
\section{符号说明}
% ==================================================================

\begin{table}[htbp]
\centering
\caption{主要符号说明}
\label{tab:symbols}
\begin{tabularx}{\textwidth}{c>{\raggedright\arraybackslash}X}
\toprule
符号 & 含义 \\
\midrule
$i$，$t$（或 $p$），$m$ & 样本序号、50个位置的序号、模态 $m\in\{T,A,V\}$（文本、声学、视觉） \\
$D_i$ & 问题一：第 $i$ 段视频的时长 \\
$[s_{ij},e_{ij})$ & 问题一：第 $j$ 个原词的起止时间 \\
$\mathcal{I}_{ip}^m$，$L_i^m$ & 问题一：位置 $p$ 的池化帧集合；模态 $m$ 的有效位置数 \\
$\ell_{imt}$，$o_{imt}$，$s_{imt}$ & 布局掩码、原生可观测掩码、人为遮蔽掩码 \\
$a_{imt}$ & 实际送入模型的可用掩码，$a=o(1-s)$ \\
$\boldsymbol{\mu}_{im}$，$\boldsymbol{d}_{im}$ & 模态 $m$ 可观测位置的均值池化与“均值+标准差”池化向量 \\
$\boldsymbol{p}_i$，$\hat c_i$，$\hat y_i$ & 三类概率、预测类别、预测强度 \\
$y_i$，$c_i$ & 真实强度（$[-3,3]$）与真实类别 \\
$r$，$\rho_{im}$ & 名义缺失比例；实际删除比例 \\
$\mathcal{K}_{PM}$ & 局部遮蔽面板：在被选模态内删除连续窗口，共63档 \\
$\mathcal{K}_{LH}$ & 整模态保留压力面板：只保留子集 $S$，其余模态整路关闭，共63档 \\
$v_i^h(S)$，$\phi_{i,m}^h$ & 只保留模态子集 $S$ 时的输出；模态 $m$ 的 Shapley 贡献，$h\in\{\mathrm{cls},\mathrm{reg}\}$ \\
$c_i^\star$，$R_k$ & 完整输入下的预测类别；局部得分最高的前 $k$ 个位置 \\
\bottomrule
\end{tabularx}
\end{table}

注：此为本文的主要符号，其他符号在正文首次出现处说明。
